\documentclass[12pt,a4paper]{article}

\usepackage[utf8]{inputenc}
\usepackage[T1]{fontenc}
\usepackage[brazil]{babel}
\usepackage{lmodern}
\usepackage{geometry}
\usepackage{booktabs}
\usepackage{array}
\usepackage{longtable}
\usepackage{amsmath}
\usepackage{graphicx}
\usepackage{hyperref}
\usepackage{xcolor}

\geometry{margin=2.5cm}
\hypersetup{
    colorlinks=true,
    linkcolor=blue,
    citecolor=blue,
    urlcolor=blue
}

\title{Infer\^encia Causal Temporal da Transfus\~ao de Hem\'acias em Pacientes Cr\'iticos de UTI}
\author{Projeto de pesquisa em Causal Machine Learning}
\date{\today}

\begin{document}

\maketitle

\begin{abstract}
Este documento resume a etapa atual de infer\^encia causal desenvolvida a partir do estudo anterior de fenotipagem fisiol\'ogica em transfus\~ao de hem\'acias na UTI. O trabalho anterior construiu coorte, ponto de decis\~ao temporal, representa\c{c}\~oes fisiol\'ogicas com MiniRocket, agrupamentos por K-means e pareamento por escore de propens\~ao. A etapa atual reaproveita esses artefatos para estimar efeitos heterog\^eneos da transfus\~ao, especialmente dentro dos grupos identificados pelo scan de benef\'icio, neutralidade e poss\'ivel malef\'icio. Os resultados indicam que a associa\c{c}\~ao entre transfus\~ao e mortalidade n\~ao \'e homog\^enea: alguns subgrupos apresentam menor mortalidade pareada entre transfundidos, enquanto outros apresentam maior mortalidade ou maior carga de suporte org\^anico.
\end{abstract}

\section{Contexto do projeto}

O projeto atual \'e uma continua\c{c}\~ao direta do artigo anterior sobre identifica\c{c}\~ao de fen\'otipos cl\'inicos e desfechos da transfus\~ao de hem\'acias em UTI baseada em trajet\'orias fisiol\'ogicas. No artigo anterior, foram utilizadas s\'eries temporais cl\'inicas pr\'e-transfusionais, representa\c{c}\~oes temporais com MiniRocket e clusteriza\c{c}\~ao por K-means. Em seguida, pacientes transfundidos foram comparados com controles por pareamento por escore de propens\~ao dentro dos fen\'otipos.

A nova etapa reorganiza esse resultado como um estudo de infer\^encia causal temporal. Em vez de refazer a constru\c{c}\~ao completa da coorte, o pipeline reaproveita os artefatos j\'a produzidos pelo projeto de fenotipagem: coorte, tabela de $t_0$, embeddings temporais, clusters, pares pareados e tabelas de desfechos. Isso evita recomputa\c{c}\~ao pesada e permite concentrar a an\'alise na pergunta causal.

\section{Pergunta causal}

A pergunta principal \'e:

\begin{quote}
Qual \'e o efeito causal da transfus\~ao de concentrado de hem\'acias sobre mortalidade e suporte org\^anico em pacientes cr\'iticos de UTI, considerando trajet\'orias fisiol\'ogicas antes do ponto de decis\~ao?
\end{quote}

A unidade de an\'alise \'e a interna\c{c}\~ao em UTI. O tratamento \'e definido como:

\[
A =
\begin{cases}
1, & \text{paciente recebeu transfus\~ao de hem\'acias;}\\
0, & \text{paciente controle n\~ao recebeu transfus\~ao.}
\end{cases}
\]

O estimando central usado nesta etapa \'e o efeito m\'edio do tratamento nos tratados, estratificado por grupo fisiol\'ogico:

\[
ATT_g = E[Y(1)-Y(0) \mid A=1, G=g],
\]

em que $G$ representa o grupo obtido por fenotipagem ou pelo scan de benef\'icio/malef\'icio, $Y(1)$ \'e o desfecho potencial sob transfus\~ao e $Y(0)$ \'e o desfecho potencial sob n\~ao transfus\~ao.

\section{Dados e artefatos utilizados}

A execu\c{c}\~ao analisada reaproveitou os resultados do pipeline anterior de fenotipagem com janela temporal de 48 horas e embedding MiniRocket. A coorte importada continha:

\begin{itemize}
    \item 40.167 interna\c{c}\~oes em UTI;
    \item 10.639 pacientes transfundidos;
    \item 29.528 controles;
    \item 10.009 pares pareados usados na estimativa pareada principal;
    \item desfecho prim\'ario: mortalidade em qualquer momento dispon\'ivel no arquivo de desfechos.
\end{itemize}

Os arquivos centrais da etapa atual s\~ao:

\begin{itemize}
    \item \texttt{outputs/causal\_inference/processed/}: coorte, tratamento, desfechos e metadados;
    \item \texttt{outputs/causal\_inference/phenotypes/}: grupos reconstru\'idos a partir dos embeddings e resultado do scan;
    \item \texttt{outputs/causal\_inference/causal/}: estimativas de efeito por grupo;
    \item \texttt{outputs/causal\_inference/counterfactual/}: modelos preditivos contrafactuais e estimativas individuais;
    \item \texttt{outputs/causal\_inference/reports/}: relat\'orios autom\'aticos.
\end{itemize}

\section{Como a infer\^encia se conecta ao artigo anterior}

O artigo anterior encontrou heterogeneidade cl\'inica: alguns fen\'otipos pareciam se associar a benef\'icio da transfus\~ao, enquanto outros sugeriam pior progn\'ostico ou maior carga de suporte. A etapa atual n\~ao substitui esse achado; ela o transforma em uma pergunta causal mais expl\'icita.

A interpreta\c{c}\~ao correta \'e:

\begin{itemize}
    \item o artigo anterior descobriu grupos fisiol\'ogicos por trajet\'orias temporais;
    \item o scan avaliou quais grupos pareciam ter benef\'icio, neutralidade ou malef\'icio;
    \item a infer\^encia atual estima o efeito da transfus\~ao dentro desses grupos;
    \item portanto, a infer\^encia est\'a sendo feita sobre subpopula\c{c}\~oes fisiologicamente definidas, n\~ao apenas sobre a coorte inteira.
\end{itemize}

\section{Resultado global pareado}

Na popula\c{c}\~ao pareada completa, a transfus\~ao esteve associada a menor mortalidade em rela\c{c}\~ao aos controles pareados. A diferen\c{c}a pareada foi de aproximadamente $-6{,}1$ pontos percentuais.

\begin{table}[ht]
\centering
\caption{Efeito pareado global da transfus\~ao.}
\begin{tabular}{lrrrr}
\toprule
Desfecho & $n$ pares & Tratados & Controles & Efeito \\
\midrule
Mortalidade & 10.009 & 0,385 & 0,446 & -0,061 \\
Ventila\c{c}\~ao, h & 10.009 & 64,3 & 32,9 & +31,4 \\
Terapia renal substitutiva & 10.009 & 0,104 & 0,065 & +0,039 \\
Noradrenalina eq. m\'axima & 10.009 & 0,265 & 0,200 & +0,065 \\
Perman\^encia na UTI, h & 10.009 & 732,8 & 421,4 & +311,4 \\
\bottomrule
\end{tabular}
\end{table}

Esse resultado global n\~ao deve ser interpretado isoladamente como evid\^encia de efeito uniforme. Apesar da redu\c{c}\~ao global de mortalidade, houve aumento importante de ventila\c{c}\~ao, terapia renal substitutiva, dose vasoativa e tempo de perman\^encia. Isso refor\c{c}a a hip\'otese de heterogeneidade: a transfus\~ao pode estar associada a menor mortalidade em parte dos pacientes, mas tamb\'em a maior carga de suporte org\^anico em outros.

\section{Resultados orientados pelo artigo: K=2}

No artigo anterior, a solu\c{c}\~ao com $K=2$ foi interpretada como dois macrofen\'otipos. Na etapa atual, esses dois grupos foram mantidos como leitura orientada pelo artigo: um grupo de benef\'icio relativo e outro de maior risco/carga de suporte.

\begin{table}[ht]
\centering
\caption{Grupos orientados pelo artigo na solu\c{c}\~ao K=2.}
\begin{tabular}{lrrrr}
\toprule
Grupo & $n$ pares & Mortalidade tratados & Mortalidade controles & Efeito \\
\midrule
Macrobenef\'icio relativo & 7.085 & 0,357 & 0,423 & -0,066 \\
Macrorrisco / alto suporte & 2.924 & 0,451 & 0,501 & -0,050 \\
\bottomrule
\end{tabular}
\end{table}

O ponto importante \'e que o segundo macrogrupo n\~ao apresentou maior mortalidade nos transfundidos em rela\c{c}\~ao aos controles pareados; ainda assim, ele apresentou carga de suporte org\^anico muito maior. Em particular, o grupo de macrorrisco/alto suporte teve aumento aproximado de 73 horas de ventila\c{c}\~ao, 9,6 pontos percentuais de terapia renal substitutiva, 0,173 na dose m\'axima equivalente de noradrenalina e 776 horas de perman\^encia em UTI.

Assim, em $K=2$, a separa\c{c}\~ao parece capturar principalmente diferen\c{c}a de gravidade e suporte org\^anico, n\~ao apenas dire\c{c}\~ao do efeito sobre mortalidade.

\section{Scan de benef\'icio, neutralidade e malef\'icio}

Para responder melhor \`a pergunta sobre subgrupos com benef\'icio ou malef\'icio, foram avaliadas solu\c{c}\~oes com $K=2,3,5,6$. A solu\c{c}\~ao $K=3$ n\~ao separou bem grupos de malef\'icio: todos os clusters ficaram classificados como benef\'icio aparente. Por isso, $K=3$ deve ser tratado como an\'alise de sensibilidade, n\~ao como solu\c{c}\~ao principal.

As solu\c{c}\~oes $K=5$ e $K=6$ foram mais informativas, pois separaram grupos com benef\'icio aparente, neutralidade e malef\'icio aparente.

\begin{table}[ht]
\centering
\caption{Efeito sobre mortalidade por grupos do scan.}
\begin{tabular}{llrrrr}
\toprule
K & Grupo & $n$ pares & Tratados & Controles & Efeito \\
\midrule
5 & Benef\'icio aparente & 6.310 & 0,361 & 0,475 & -0,114 \\
5 & Malef\'icio aparente & 1.813 & 0,461 & 0,391 & +0,069 \\
5 & Neutro aparente & 1.886 & 0,389 & 0,399 & -0,011 \\
\midrule
6 & Benef\'icio aparente & 5.340 & 0,321 & 0,461 & -0,140 \\
6 & Malef\'icio aparente & 2.836 & 0,499 & 0,447 & +0,052 \\
6 & Neutro aparente & 1.833 & 0,393 & 0,399 & -0,005 \\
\bottomrule
\end{tabular}
\end{table}

Esses achados sustentam a hip\'otese principal do projeto: a transfus\~ao n\~ao parece ter efeito homog\^eneo. Na solu\c{c}\~ao $K=5$, o grupo de benef\'icio aparente teve redu\c{c}\~ao de aproximadamente 11,4 pontos percentuais na mortalidade, enquanto o grupo de malef\'icio aparente teve aumento de aproximadamente 6,9 pontos percentuais. Na solu\c{c}\~ao $K=6$, o grupo de benef\'icio aparente teve redu\c{c}\~ao de aproximadamente 14,0 pontos percentuais, enquanto o grupo de malef\'icio aparente teve aumento de aproximadamente 5,2 pontos percentuais.

\section{Suporte org\^anico nos grupos K=5 e K=6}

Al\'em da mortalidade, os grupos tamb\'em diferiram quanto a suporte org\^anico. Na solu\c{c}\~ao $K=6$, o grupo de malef\'icio aparente combinou maior mortalidade com aumento de ventila\c{c}\~ao, terapia renal substitutiva, dose vasoativa m\'axima e perman\^encia em UTI.

\begin{table}[ht]
\centering
\caption{Desfechos secund\'arios principais na solu\c{c}\~ao K=6.}
\begin{tabular}{llrr}
\toprule
Grupo & Desfecho & Efeito & Interpreta\c{c}\~ao \\
\midrule
Benef\'icio aparente & Mortalidade & -0,140 & menor mortalidade \\
Benef\'icio aparente & Ventila\c{c}\~ao, h & +28,7 & maior suporte \\
Benef\'icio aparente & Perman\^encia UTI, h & +265,9 & maior tempo de UTI \\
\midrule
Malef\'icio aparente & Mortalidade & +0,052 & maior mortalidade \\
Malef\'icio aparente & Ventila\c{c}\~ao, h & +40,8 & maior suporte \\
Malef\'icio aparente & Terapia renal substitutiva & +0,097 & maior suporte renal \\
Malef\'icio aparente & Noradrenalina eq. m\'axima & +0,089 & maior intensidade vasoativa \\
Malef\'icio aparente & Perman\^encia UTI, h & +447,1 & maior tempo de UTI \\
\midrule
Neutro aparente & Mortalidade & -0,005 & efeito pequeno \\
Neutro aparente & Ventila\c{c}\~ao, h & +24,6 & maior suporte \\
\bottomrule
\end{tabular}
\end{table}

Esse padr\~ao sugere que a mortalidade isolada n\~ao captura toda a heterogeneidade cl\'inica. Mesmo nos grupos classificados como benef\'icio aparente, os pacientes transfundidos podem apresentar maior tempo de suporte org\^anico, possivelmente refletindo maior sobrevida com maior tempo em risco, maior gravidade residual ou confus\~ao n\~ao completamente removida.

\section{Modelos contrafactuais e CATE/ITE}

Tamb\'em foi implementado um modelo contrafactual inicial com dois modelos preditivos: um para estimar $Y(1)$ e outro para estimar $Y(0)$. Esse modelo gerou risco esperado sob transfus\~ao, risco esperado sob n\~ao transfus\~ao e efeito individual estimado:

\[
ITE_i = \hat{Y}_i(1) - \hat{Y}_i(0).
\]

O modelo preditivo baseado em gradient boosting apresentou desempenho razo\'avel para mortalidade:

\begin{table}[ht]
\centering
\caption{Desempenho preditivo do modelo contrafactual inicial.}
\begin{tabular}{lrr}
\toprule
M\'etrica & Conjunto completo & Teste \\
\midrule
AUC & 0,869 & 0,849 \\
Average precision & 0,780 & 0,732 \\
Brier score & 0,135 & 0,145 \\
F1 & 0,673 & 0,637 \\
\bottomrule
\end{tabular}
\end{table}

Apesar do bom desempenho preditivo, essas estimativas individuais ainda devem ser interpretadas como explorat\'orias. Um modelo pode prever bem o desfecho observado sem necessariamente identificar corretamente o desfecho contrafactual n\~ao observado.

Tamb\'em foi testado um modelo de efeito heterog\^eneo no estilo causal forest. Como bibliotecas especializadas como \texttt{econml} ou \texttt{causalml} n\~ao estavam dispon\'iveis, foi utilizado um fallback com T-learner baseado em gradient boosting. Esse modelo n\~ao reproduziu fortemente os sinais encontrados no scan pareado, sugerindo que os efeitos individuais s\~ao sens\'iveis ao m\'etodo e ainda precisam de valida\c{c}\~ao.

\section{Interpreta\c{c}\~ao cient\'ifica}

Os resultados atuais apontam tr\^es conclus\~oes principais.

Primeiro, a an\'alise global sugere associa\c{c}\~ao entre transfus\~ao e menor mortalidade na popula\c{c}\~ao pareada. No entanto, essa estimativa global esconde heterogeneidade relevante.

Segundo, os grupos derivados do scan em $K=5$ e $K=6$ indicam subpopula\c{c}\~oes com dire\c{c}\~oes opostas de efeito. Isso \'e compat\'ivel com a hip\'otese de que a decis\~ao transfusional deve depender da trajet\'oria fisiol\'ogica pr\'e-interven\c{c}\~ao, e n\~ao apenas de limiares fixos de hemoglobina.

Terceiro, a solu\c{c}\~ao $K=2$ do artigo continua relevante como macrofenotipagem, mas ela parece resumir uma separa\c{c}\~ao cl\'inica ampla entre menor e maior carga de suporte. Para infer\^encia sobre benef\'icio e malef\'icio, as solu\c{c}\~oes $K=5$ e $K=6$ s\~ao mais informativas.

\section{Limita\c{c}\~oes}

As estimativas s\~ao observacionais. Portanto, dependem de ignorabilidade condicional, positividade, consist\^encia, qualidade do pareamento e aus\^encia de confundimento residual. Al\'em disso, os grupos de benef\'icio e malef\'icio foram obtidos por scan explorat\'orio; por isso, devem ser tratados como geradores de hip\'otese, n\~ao como prova confirmat\'oria definitiva.

Outro ponto importante \'e que os metadados da execu\c{c}\~ao importada indicavam \texttt{keep\_post\_t0\_features=true}. Assim, os modelos preditivos e contrafactuais devem ser considerados explorat\'orios at\'e uma execu\c{c}\~ao estritamente pr\'e-$t_0$. A vers\~ao confirmat\'oria deve remover qualquer vari\'avel derivada ap\'os $t_0$ antes de treinar modelos de propens\~ao, CATE ou contrafactual.

\section{Conclus\~ao}

O projeto evoluiu de uma an\'alise de fenotipagem fisiol\'ogica para um pipeline de infer\^encia causal temporal. O principal achado \'e que a transfus\~ao de hem\'acias parece ter associa\c{c}\~oes heterog\^eneas com mortalidade e suporte org\^anico. Em particular, os grupos $K=5$ e $K=6$ revelaram subgrupos com benef\'icio aparente, neutralidade e malef\'icio aparente.

O pr\'oximo passo metodol\'ogico \'e consolidar uma execu\c{c}\~ao confirmat\'oria com vari\'aveis estritamente pr\'e-$t_0$, diagn\'ostico de overlap, balanceamento por grupo, bootstrap, sensibilidade a caliper e compara\c{c}\~ao entre janela de 24 e 48 horas. Essa etapa permitir\'a transformar os achados explorat\'orios em uma an\'alise causal mais robusta.

\end{document}
